% =============================================================================
% CARTOGRASS 2026 poster -- TEXTS: materials and methods.
% -----------------------------------------------------------------------------
% Order in the Materials-and-methods row (2026-10-19): left half, the study
% sites (one strip map), what goes into MAREA and the independent checks,
% then one real pixel ("A pixel is a tide gauge"); right half, the method as
% six cards that follow one imaginary pixel (figures/cg/method_flow,
% tools/cg/fig_method_flow.py). Reference data are described here; their
% results go in text_results.tex (placeholders).
% Audience pass (2026-10-19, CARTOGRASS: seagrass ecologists, cartographers,
% managers): plain words, one word for one thing. "Delay", not "lag", in every
% visible text; "section", not "band", for the 3-km delay units (nobody reads
% a spectral band); no CDSE, SCL classes or NDWI without its meaning (the flow
% chart says "water index from green and near-infrared light"); "never used to
% make the maps", not "never used in the fit".
% Round 2 (2026-10-19, review of the method panels): the data note says the
% practical gain for managers (no field data); the checks name "each site"
% (the Wadden site is a tidal basin, not an estuary); the pixel panel's note
% points at step 4 of the flow chart; one look rule for the whole sheet.
% Round 3 (2026-10-19): the inputs carry the chart's INPUT tags, the tide's
% source is named per site ("at four sites"), "no field surveys"; the
% fallback steps follow the chart's new card wordings.
% Both lists are body text (\normalsize, 27.9 pt), set as run-in paragraphs.
% =============================================================================
\ifdefined\cgMethodsTextLoaded\endinput\fi
\def\cgMethodsTextLoaded{}

% A bare itemize prints NO bullet in this framework (beamer's item template is
% lost once enumitem owns itemize), so a list passes the theme's teal square.
\providecommand{\cgItemLabel}{{\color{PTeal}\usebeamertemplate{itemize item}}}

\newcommand{\cgMethodsHeading}{Materials and methods}

% --- What goes in ------------------------------------------------- 30 words
% Facts behind the words (unchanged from 2026-10-16): Sentinel-2 L2A, 10 m,
% 2023-2025 (Sentinel-2A, 2B, 2C), processed with openEO [3] on the Copernicus
% Data Space Ecosystem ("the Copernicus cloud"); a look is a clear pixel by
% the Sentinel-2 scene classification (classes 4-7 kept) of a scene with
% little cloud (at most 5 % cloudy over the frame or 10 % over the transition
% zone). Tide at the mouth: EOT20 [4], blended with the outer IOC gauges where
% they exist (Ferrol fer1; Westerschelde vlis + brsk; Wadden ters; Ems bork;
% Villaviciosa EOT20 alone), hence "where they exist".
% Its note says what the brief says of the method, in the words that matter
% to monitoring and management: MAREA needs no field survey. Nothing is
% trained on labelled data; every delay and elevation comes from each site's
% own images and tide (the checks below never enter: their note says so).
% Round 3 (2026-10-19): "no field surveys", not "no field data": the tide
% gauges at the mouth are in-situ data, and they are an input.
% Round 4 (2026-10-19): "tide gauges at the mouth", not "outside the
% estuary": Ferrol 1 stands in the outer ria, Vlissingen + Breskens and
% Borkum Fischerbalje at the estuary mouths; only Terschelling Noordzee is
% truly outside. "At the mouth" is also INPUT 2's own label.
% The two inputs carry the flow chart's navy tabs (\cgInputTag, cg_macros.tex:
% the chart's cards 1 and 2 are "INPUT 1 · IMAGES" and "INPUT 2 · TIDE AT THE
% MOUTH"), so the chart's tide input finds its source here. Gauges: the
% production boundary blends EOT20 with outer gauges at four sites (Ferrol,
% Westerschelde, Wadden, Ems); at Villaviciosa it is EOT20 alone
% (experiments/benchmark/config.py BOUNDARY), hence "at four sites".
\newcommand{\cgDataTitle}{Input data}
\newcommand{\cgDataNote}{(all MAREA needs: no field surveys)}
\newcommand{\cgDataBody}{\cgInputTag{1}\ Sentinel-2 images, 10\,m, 2023--2025,
  as openEO data cubes in the Copernicus cloud~[3].\quad \cgInputTag{2}\ Tide at
  the mouth: global tide model EOT20~[4] and, at four sites, tide gauges at
  the mouth.}

% --- Independent checks (materials) ------------------------------- 30 words
% Never used to make the maps (hence "independent"); the outer gauges that
% feed the tide at the mouth are inputs and are named above. The numbered
% discs (\cgSiteNum, cg_macros.tex) are the study-site map's. No point count
% for the RTK survey here: it holds 361 points, 270 are development points and
% 194 of those fall on mapped pixels (the number scored, in the pilot box).
% Glosses for the audience: RTK-GNSS = a precise GPS survey; PNOA = the
% Spanish national airborne LiDAR; Vaklodingen = Rijkswaterstaat's depth
% surveys of the Dutch coast.
% Round 4 (2026-10-19): the gauges named here are the 13 headline judges of
% the frozen validation (experiments/benchmark/judges.HEADLINE;
% docs/validation_v2_spec.md: Borssele and Eemshaven haven are excluded for
% coverage, fer4 for standing less than 3 km from the boundary gauge Ferrol 1),
% the only triangles the site strip draws. Ferrol has one, Ferrol 2 in the
% inner port: "a tide gauge in the inner ria".
\newcommand{\cgSitesTitle}{Study sites and data}
\newcommand{\cgRefTitle}{Independent checks}
\newcommand{\cgRefNote}{(never used to make the maps)}
\newcommand{\cgSitesBody}{%
  \cgSiteNum{1}\,precise GPS survey (RTK-GNSS) and PNOA airborne LiDAR\quad
  \cgSiteNum{2}\,a tide gauge in the inner ria\quad
  \cgSiteNum{3}\,\cgSiteNum{4}\,\cgSiteNum{5}\,Vaklodingen depth surveys and
  tide gauges inside each site}

% --- How MAREA works (the flow chart) -------------------------------------
% Title and note of the right-hand panel. The note says that the pixel and
% its looks are simulated (fig_method_flow.py: planted delay 90 min, recovered
% 90 min; the minutes shown are examples, not results).
\newcommand{\cgStepsTitle}{How MAREA works: follow one pixel}
\newcommand{\cgStepsNote}{Schematic: one imaginary pixel}
% Fallback only: set in the figure's place while figures/cg/method_flow is
% missing (blk_mm_method.tex). Same six steps as the figure; step 3 also keeps
% the mouth and monotone rule (m2a_rasch: tau = 0 at the mouth, monotone=True),
% which the figure no longer draws since 2026-10-20.
\newcommand{\cgStepsBody}{%
  \step{1}{\term{Seen many times.} Each cloud-free Sentinel-2 view of the
    pixel says wet or dry (a water index from green and near-infrared light).}
  \step{2}{\term{Its tide arrives late.} The tide at the mouth (global tide
    model and gauges) reaches its 3-km section some minutes later: the pixel
    is wet when that water stands above its ground. MAREA shifts that tide
    in time, not in height.}
  \step{3}{\term{Its section's delay.} One delay for the whole section: the
    one that best splits the wet and dry looks of its pixels, judged
    together; zero at the mouth, it can only grow inland.}
  \step{4}{\term{Its elevation.} The midpoint of its flooding curve on the
    corrected tide; the width of the curve is the uncertainty. A pixel without
    both wet and dry looks is left empty.}
  \step{5}{\term{Its emerged time.} The share of time the corrected tide stays
    below its ground, at all times of the tide record, not only at image
    times.}
  \step{6}{\term{All pixels: the maps.} Elevation with its uncertainty, and
    emerged time: the base layer for the potential habitat of
    \emph{Zostera noltei}.}}

% --- A pixel is a tide gauge -----------------------------------------------
% One real 10 m pixel near the mouth of the Ria de Villaviciosa
% (tools/cg/fig_pixel_gauge.py): its 95 looks in MAREA's own 2023-2025 record
% (97 scenes kept by the record's scene rule, the pixel's own cloud class
% clear: the flow chart's "1 look = 1 cloud-free view"); 2.9 km from the mouth
% along the water its delay is negligible, so its water is the tide at the
% mouth (the figure says "3 km from the mouth: delay tau ~ 0"; its tide axis
% says "mouth tide", the flow chart's h). The note points at the flow chart's
% step 4, which this panel shows on real data. Round 3: "the method's step",
% because the left half is read first, before the chart that numbers the
% steps, and a bare teal 4 sat a few cm from the sand site disc 4.
\newcommand{\cgPixTitle}{A pixel is a tide gauge}
\newcommand{\cgPixNote}{the method's step \cgStepNum{4} on one real pixel, Ría
  de Villaviciosa, 2023--2025}
